\documentclass[10pt,a4paper]{amsart}
\usepackage[margin=19mm]{geometry}
\usepackage{amsmath,amssymb,mathtools}
\usepackage[scr=rsfs,cal=euler]{mathalfa}
\usepackage{xfrac}
\usepackage[unicode,hidelinks,bookmarks=false]{hyperref}
\newcommand{\ie}{i.e.\ }
\newcommand{\cf}{cf.\ }
\newcommand{\resp}{respectively\ }
\newcommand{\Subsection}{\subsection}
\providecommand{\nobreakdash}{\nobreak\hbox{-}}
\setlength{\emergencystretch}{2em}
\multlinegap0pt
\usepackage{amssymb}
\usepackage{float}
\usepackage{enumerate}
\usepackage{xcolor}
\newtheorem{lemma}{Lemma}
\title{Finite groups with mostly involuted cyclic subgroups}
\author{Vaibhav Chhajer, Palash Sharma}
\date{Source: arXiv:2508.00681v3 (CC BY 4.0)}
\begin{document}
\maketitle

\noindent\emph{Source and licence.} Finite groups with mostly involuted cyclic subgroups. Version arXiv:2508.00681v3 (CC BY 4.0), distributed by arXiv under the Creative Commons Attribution 4.0 International licence, http://creativecommons.org/licenses/by/4.0/, as recorded in the arXiv licence metadata for this exact version and shown on the abstract page https://arxiv.org/abs/2508.00681v3. Full licence notice: \texttt{licenses/arxiv-2508.00681-CC-BY-4.0.txt}.\par\smallskip

\noindent\textit{Editorial scope.} Counted unit: the article's lemma beta (source0.tex lines 335--339). The article's complete original proof of it is delivered with it (source0.tex lines 341--364, one \texttt{proof} environment of 24 lines). No mathematics is written, repaired or re-derived: the statement and the proof are the article's own lines, in the article's own notation, and no retained line is substituted or re-flowed.  No further source-local context is needed: every cross-reference of every retained line resolves inside the two delivered spans.  Nothing was omitted from any retained span, so the package carries no omission record.  No retained line contains a citation, so the wrapper carries no bibliography and no bibliography entry.  No retained line contains a figure environment, an image-inclusion command or drawing code, so no image is delivered and the package declares no asset dependency.  Editorial formatting only, in addition: an AMS \texttt{amsart} preamble that restates the article's own declarations and macros used by the retained lines, namely the environments \texttt{lemma} on separate counters, together with the standard packages those lines need.  The retained environments print in this document as Lemma 1 in order of appearance, whereas the article prints the same items with the numbering of its own full document.  The complete native source of the article as distributed in the arXiv e-print for this exact version is bundled unchanged as \texttt{sources/182563/source0.tex}.
\begin{lemma} \label{basic property of beta}
    Let $G$ and $H$ be finite groups. Then 
    $$\beta(G\times H)=\beta(G)\cdot\beta(H)\ 
    \text{if and only if } \operatorname{gcd}(\operatorname{exp}(G),\operatorname{exp}(H)) \in \{1,2\}.$$
\end{lemma}

\begin{proof}
    Clearly, $i(G\times H) = i(G)\cdot i(H)$ holds unconditionally. Therefore, $\beta(G\times H)=\beta(G)\cdot\beta(H)$ if and only if $c(G\times H)=c(G)\cdot c(H)$. 
    
    Recall that for any finite group $K$, $c(K) = \sum_{x \in K} \frac{1}{\phi(o(x))}$. Thus, 
    $$c(G\times H) = \sum_{(x,y)\in G\times H} \frac{1}{\phi(\operatorname{lcm}(o(x),o(y)))}.$$
    Using the property of $\phi$ that $\phi(\operatorname{lcm}(m,n))\cdot \phi(\operatorname{gcd}(m,n)) =\phi(m)\cdot\phi(n)$ for all $m,n \in \mathbb{N}$, we obtain
    $$c(G\times H) = \sum_{(x,y)\in G\times H} \frac{\phi(\operatorname{gcd}(o(x),o(y)))}{\phi(o(x))\cdot\phi(o(y))}.$$
    Since $\phi(k) \ge 1$ for all $k \in \mathbb{N}$, we always have $c(G\times H) \ge c(G)\cdot c(H)$, with equality holding if and only if $\phi(\operatorname{gcd}(o(x),o(y))) = 1$ for all $x \in G$ and $y \in H$, that is, if and only if $ \operatorname{gcd}(o(x),o(y)) \in \{1,2\}$ for all $x \in G$ and $y \in H$. Therefore, it suffices to show that $ \operatorname{gcd}(o(x),o(y)) \in \{1,2\}$ for all $x \in G$ and $y \in H$ holds if and only if $\operatorname{gcd}(\operatorname{exp}(G),\operatorname{exp}(H)) \in \{1,2\}$. 
    
    If $\operatorname{gcd}(\operatorname{exp}(G),\operatorname{exp}(H)) \in \{1,2\}$, then since $o(x)$ divides $\operatorname{exp}(G)$ and $o(y)$ divides $\operatorname{exp}(H)$, it immediately follows that $\operatorname{gcd}(o(x),o(y))$ divides $\operatorname{gcd}(\operatorname{exp}(G),\operatorname{exp}(H))$, and thus $\operatorname{gcd}(o(x),o(y))$ $\in \{1,2\}$.
    
    Conversely, suppose that $\operatorname{gcd}(o(x), o(y)) \in \{1, 2\}$ for all $x \in G$ and $y \in H$, but, for the sake of contradiction, that $\operatorname{gcd}(\operatorname{exp}(G), \operatorname{exp}(H)) \notin \{1, 2\}$. Then we have the following two cases:
    
    \textbf{Case 1:} Suppose $\operatorname{gcd}(\operatorname{exp}(G), \operatorname{exp}(H)) = 2^n r$, where $n \geq 0$ and $r \geq 3$ is an odd integer.
Then, an odd prime $p$ divides $\operatorname{gcd}(\operatorname{exp}(G), \operatorname{exp}(H))$ and so $p \mid |G|$ and $p \mid |H|$. Thus, there exist $a \in G$ and $b \in H$ such that $o(a) = p$ and $o(b) = p$. Consequently, $\operatorname{gcd}(o(a), o(b)) = p \notin \{1, 2\}$, a contradiction.

    
  \textbf{Case 2:} Suppose $\operatorname{gcd}(\operatorname{exp}(G), \operatorname{exp}(H)) = 2^n$, where $n \geq 2$. Then, $4 \mid \operatorname{exp}(G)$ and $4 \mid \operatorname{exp}(H)$. 
If it were true that $o(x) \leq 2$ for all $x \in G$, this would imply $\operatorname{exp}(G) \leq 2$, which contradicts $4 \mid \operatorname{exp}(G)$. Therefore, there exists $a \in G$ such that $o(a)=4$. Similarly, there exists $b \in H$ such that $o(b)=4$.
Consequently, $\operatorname{gcd}(o(a), o(b)) = 4 \notin \{1, 2\}$, a contradiction.

    
    Since both cases lead to a contradiction, we conclude that $\operatorname{gcd}(\operatorname{exp}(G), \operatorname{exp}(H)) \in \{1, 2\}$.
\end{proof}

\end{document}
